\section{Doctrine of the Soul}

\label{sec:DoctrineSoul}

In the Medieval period, the doctrine of the soul reached its deepest understanding in the thought of \textbf{Thomas Aquinas}. This doctrine was adopted, virtually \emph{in toto}, from Aristotle. Now, \textbf{Aristotle} was the heir to the totality of Greek philosophy and its mysteries. Even if much of that earlier work has now been lost except in fragments, it was familiar to Aristotle, who absorbed and organized it. Furthermore, as Guenon takes pain to demonstrate on numerous occasions, Aristotle is in deep agreement with the metaphysicians of India. As we have recently pointed out, this may be the result of a direct historical influence.

Bear in mind that the ultimate and primary goal of the Medieval Christian is to save his soul. But we see that the understanding of that soul is purely pagan in its origins. And this understanding insinuates itself at the heart of Medieval doctrine.

There is no point to repeat what was written in \textit{The Intellectual Soul}\footnote{See Section \ref{sec:IntellectualSoul} in this book.}. There are a few important implications of this doctrine.

\begin{itemize}
\item The soul forms the body, not the other way around as most contemporary Westerners seem to believe (if they even consider the question). 
\item The body and the components of the soul form a unity. 
\item The natural order is for the body to be subservient to the soul, and each layer of the soul to be subservient to the one above. 
\item Unlike the view of some gnostic doctrines, and even Plato, the soul is not ``imprisoned" in the body. There is a struggle between the body and soul to enforce the natural order between them. However, this cannot be understood as being ``anti-life". 
\item The intellectual soul reveals ``man as the one who possesses the power of knowing with respect to the visible world." (TB 5:6) This knowing is transcendental to the world. Guenon writes, in reference to Aristotle: ``Pure intellect is of a transcendent order, knowledge of universal principles is its proper object." 
\item The body is the object of the Will. Man is a subject not just in knowing, but also in acting. ``The structure of this body is such that it permits him to be the author of all human activity." (TB 7:2) 
\item The human person is never experienced as a phenomenon, but is revealed indirectly through his activity. Hence, the person is transcendent. Again, Guenon relies on Aristotle when he asserts, ``The personality (\emph{Atma}) is not manifested", since it is the principle of manifestation. 
\item For Aristotle, the principle of all action must be actionless. That is, the unmoved mover is not an object in the action he initiates. The person, then, is that unmoved mover, the microcosmic version of the macrocosm. 
\end{itemize}


One step to fully recover Tradition, then, is to use Guenon's Man and his Becoming to supplement and augment the Western understanding. However, it is clear that for the Westerner, this explanation is preferable to slogging through such a dense work, replete with Sanskrit words. Unfortunately, and perhaps necessarily, there are many Westerners who look to the East, not realizing what lies closer to home. The real benefit would come from those who understand the Eastern teachings and then re-express them in a Western context. Have we mentioned Aristotle and Aquinas as examples of how this can be done?

For example, one could take what Guenon writes about the envelopes of the Self and apply it to the Western understanding of the soul. Obviously, I am assuming a familiarity with the Western tradition, which is far from common knowledge. The teachings of Purusha and Prakriti from the Sankhya school of Vedanta have direct analogs with the hylomorphism of Aristotle and Aquinas. Purusha is the form or idea that acts on Prakriti, or prime matter.

Matter for the Medievals meant something quite different from what is meant today. Prime matter, or Prakriti, has no existence on its own, but is pure potentiality. Hence, things are not ``made" out of matter. Just the opposite, a thing is manifested when its idea informs the potential of prime matter, making it actual. Some may object that this is unscientific. Perhaps for the outmoded mechanistic view of matter, but it is compatible with quantum theory. This theory, following Schrodinger, holds that matter is really energy waves in an indeterminate state. Only a conscious observer, or as we prefer, a conscious actor, collapses the wave, making a determinate thing.

Since the soul is not exhausted in its bodily functions, it has two modes of existence, a corporeal one and a purely spiritual one. Hence, we can speak of an \emph{anima spiritualis}, or spirit-soul. But this is a topic for another time.

The homework is to make the interior shift in worldview to understand oneself and the world in this way. Through self-knowledge and self-observation, learn to distinguish in consciousness the various sheaths of the soul. Where is an urge arising from, what creates a certain emotion, what is the source of a thought? Discover the ``person", the unmoved mover, that is not directly experienced, yet is part of every experience. Change your attitude about matter, which has no real existence. Experience everything as the result of the ideas that inform things and events. Matter itself is empty. What we experience as resistance from matter is really a Privation\footnote{\url{https://www.gornahoor.net/?p=1815}}, something that our Will is not yet strong enough to overcome. The spirit-soul has great power, or siddhis, \emph{in potentia}; that can be harnessed.

References:

Rene Guenon, \textit{Man and his Becoming}

Pope JP II, \textit{Theology of the Body (TB)}

Alois Wiesinger, OCSO, \textit{Occult Phenomena}



\flrightit{Posted on 2012-01-20 by Cologero }

\begin{center}* * *\end{center}

\begin{footnotesize}\begin{sffamily}



\texttt{Charlotte Cowell on 2012-01-20 at 15:25 said: }

when do you think the `way of Fool' came into play, and with whom?


\hfill

\texttt{Charlotte Cowell on 2012-01-22 at 15:47 said: }

While I remember – seeing as a defence of scholasticism is in order and also that I kept thinking over the past few weeks I must somehow get us onto Lazarus – here is an extract from VT's wonderful book, Lazarus, come forth!,that I thought was very pertinent:

\url{http://alchemical-weddings.com/alchemical-weddings/lazurus-come-forth}


\hfill

\texttt{Logres on 2012-01-22 at 20:53 said: }

Sir, was the blank spot at the end meant to be a spur to personal meditation? Or could you supply the ending, if there is one? Helpful article – I find it like pulling teeth to try to discern the parts of the soul – very confusing at first. Like unsnarling a knot.


\hfill

\texttt{Cologero on 2020-01-20 at 10:55 said: }

Theosophical inspired and other New Age systems usually have the false, if not absurd, goal of ``contacting the Higher Self". The Higher Self can never be an ``object" in the world; hence it makes no sense to try to ``contact" it.


\end{sffamily}\end{footnotesize}
